% TOP_PROBLEM: {"id": 4700004, "title": "Low-degree classical Liénard systems", "category_id": 11, "category": {"id": 11, "name": "dynamical_systems", "display_name": "Dynamical Systems", "description": "Problems about long-term behavior of deterministic systems, Hamiltonian dynamics, and periodic orbits.", "slug": "dynamical-systems", "order_index": 11, "created_at": "2026-07-31T15:26:25.671Z"}, "status": "partially_solved", "problem_number": "AMR-046-0004", "set_id": 13, "statusQualification": "The numerical equalities are established; the source asks for proofs using Dulac functions."}
% FIRST_POSTED: 2026-10-02
% ENTRY_KIND: statement_only
% UnsolvedMath ID 4700004; source catalogue code AMR-046-0004
% !TeX program = lualatex
% Definition and sources only; this file is not an LLM solution attempt.
\documentclass[11pt,a4paper]{article}
\usepackage[margin=25mm]{geometry}
\usepackage{fontspec}
\setmainfont{lmroman10-regular.otf}[BoldFont=lmroman10-bold.otf,
 ItalicFont=lmroman10-italic.otf,BoldItalicFont=lmroman10-bolditalic.otf]
\usepackage{amsmath,amssymb,mathtools}
\usepackage{unicode-math}
\setmathfont{latinmodern-math.otf}
\AtBeginDocument{\let\setminus\smallsetminus}
\usepackage{xurl,hyperref}
\hypersetup{colorlinks=true,urlcolor=blue}
\setcounter{secnumdepth}{0}
\setlength{\parindent}{0pt}
\setlength{\parskip}{5pt}
\setlength{\emergencystretch}{3em}
\begin{document}
\section*{Low-degree classical Liénard systems}\label{4700004}
{\small UnsolvedMath ID 4700004 · AMR-046-0004 · Dynamical Systems · AMR Open Problem Lists}
\subsection{Definitions and mathematical statement}
Let $F\in\mathbb R[x]$ have degree $n$ and normalize $F(0)=0$; subtracting its
constant term merely translates the $y$ coordinate. Consider the classical
Liénard system
\[
 \dot x=y-F(x),\qquad \dot y=-x.
\]
Write $\operatorname{Lie}(n)$ for the supremum of the number of its isolated
nonconstant periodic orbits over all such $F$. The equalities
$\operatorname{Lie}(3)=\operatorname{Lie}(4)=1$ are known. The question asks
for proofs of these two equalities using the Bendixson–Dulac method.

For a vector field $X=(P,Q)$, a Dulac multiplier is a differentiable scalar
function $B$ on an appropriate region such that
\[
 \operatorname{div}(BX)=\partial_x(BP)+\partial_y(BQ)
\]
has a controlled sign. On a simply connected region, a suitable strict-sign
condition rules out periodic orbits by Green's theorem. For uniqueness one
may partition the plane along a zero or singular set of $B$, allowing a
region with one hole and controlling any closed trajectories on its boundary.

The task is to construct suitable multipliers, or equivalent auxiliary
functions, and verify the sign and topology conditions uniformly for every
cubic and quartic $F$. Cases where the origin is a center must be handled:
their non-isolated periodic orbits do not count as limit cycles. This is a
request for a particular kind of proof of known numerical bounds, together
with the usual examples showing the bound one is attained.
\subsection{Short English statement}
Find Bendixson–Dulac proofs that cubic and quartic classical Liénard systems have at most one limit cycle.
\subsection{Sources}
\begin{itemize}
\item[S1] ULAM.AI and the UnsolvedMath Contributors, supplied problems.json, record 4700004 (AMR-046-0004), AMR Open Problem Lists. Catalogue identity and source transcription; CC BY 4.0.
\url{https://huggingface.co/datasets/ulamai/UnsolvedMath}.
\item[S2] Gasull - Some open problems in low dimensional dynamical systems (2020), Problem environment 4: Low-degree classical Liénard systems. Original source indexed by AMR-046-0004.
\url{https://arxiv.org/abs/2012.02524}.
\end{itemize}
\end{document}
